\textbf{Phần 1. Trắc nghiệm lựa chọn}
\setcounterpageref{ex}{0}

\begin{ex} %[2D1V01]
Tìm một số có hai chữ số, biết chữ số hàng chục là $6$ và chữ số hàng đơn vị lớn hơn chữ số hàng chục $3$ đơn vị.
\choice
{$96$}
{$63$}
{\True $69$}
{$68$}
\loigiai{
Chữ số hàng chục là $6$. Chữ số hàng đơn vị là $6 + 3 = 9$. Vậy số đó là $69$.
\textbf{Đáp án của câu hỏi là C}
}
\end{ex}

\begin{ex} %[2D1V01]
Từ ba chữ số $2, 0, 7$, lập được tất cả bao nhiêu số có hai chữ số khác nhau?
\choice
{\True $4$ số}
{$6$ số}
{$5$ số}
{$3$ số}
\loigiai{
Chữ số hàng chục có thể chọn là 2 hoặc 7 (không thể chọn 0).\n- Hàng chục là 2: lập được $20, 27$.\n- Hàng chục là 7: lập được $70, 72$.\nCó tất cả $4$ số.
\textbf{Đáp án của câu hỏi là A}
}
\end{ex}

\begin{ex} %[2D1V01]
Một số có hai chữ số mà tổng hai chữ số bằng $12$, hiệu giữa chữ số hàng chục và chữ số hàng đơn vị bằng $0$. Số đó là
\choice
{$57$}
{$48$}
{$75$}
{\True $66$}
\loigiai{
Hiệu hai chữ số bằng 0 nghĩa là chữ số hàng chục bằng chữ số hàng đơn vị.\nTổng hai chữ số bằng $12$, mà $6 + 6 = 12$. Vậy số đó là $66$.
\textbf{Đáp án của câu hỏi là D}
}
\end{ex}

\begin{ex} %[2D1V01]
Cho các số: $25, 62, 52, 19, 70$. Khi sắp xếp các số này theo thứ tự từ lớn đến bé, số đứng ở vị trí thứ ba là
\choice
{$62$}
{\True $52$}
{$25$}
{$70$}
\loigiai{
Sắp xếp từ lớn đến bé: $70 > 62 > 52 > 25 > 19$. Số đứng thứ ba là $52$.
\textbf{Đáp án của câu hỏi là B}
}
\end{ex}

\begin{ex} %[2D1V01]
Số tròn chục liền trước của số $80$ là
\choice
{$79$}
{$81$}
{\True $70$}
{$60$}
\loigiai{
Các số tròn chục liên tiếp là $60, 70, 80, 90$. Số tròn chục liền trước 80 là $70$.
\textbf{Đáp án của câu hỏi là C}
}
\end{ex}

\begin{ex} %[2D1V01]
Có bao nhiêu số có hai chữ số mà chữ số tận cùng là $5$?
\choice
{\True $9$ số}
{$10$ số}
{$8$ số}
{$5$ số}
\loigiai{
Các số đó là: $15, 25, 35, 45, 55, 65, 75, 85, 95$. Có tất cả $9$ số.
\textbf{Đáp án của câu hỏi là A}
}
\end{ex}

\begin{ex} %[2D1V01]
Tìm số nhỏ nhất có hai chữ số mà chữ số hàng chục hơn chữ số hàng đơn vị $5$ đơn vị.
\choice
{$61$}
{$72$}
{$49$}
{\True $50$}
\loigiai{
Để số đó nhỏ nhất, chữ số hàng đơn vị phải nhỏ nhất, tức là bằng $0$.\nKhi đó chữ số hàng chục là $0 + 5 = 5$. Vậy số nhỏ nhất là $50$.
\textbf{Đáp án của câu hỏi là D}
}
\end{ex}

\begin{ex} %[2D1V01]
Cho dãy số: $3, 7, 11, 15, \dots, 27$. Hai số thích hợp cần điền vào chỗ chấm là
\choice
{$18, 22$}
{\True $19, 23$}
{$19, 24$}
{$17, 21$}
\loigiai{
Quy luật dãy số: mỗi số đứng sau bằng số đứng trước cộng thêm $4$.\n$15 + 4 = 19$; $19 + 4 = 23$; $23 + 4 = 27$. Vậy điền $19, 23$.
\textbf{Đáp án của câu hỏi là B}
}
\end{ex}

\begin{ex} %[2D1V01]
Một rổ táo có nhiều hơn $35$ quả nhưng ít hơn $42$ quả. Nếu đếm theo chục thì được $4$ chục. Hỏi trong rổ có bao nhiêu quả táo?
\choice
{$36$ quả táo}
{$38$ quả táo}
{\True $40$ quả táo}
{$41$ quả táo}
\loigiai{
$4$ chục bằng $40$. Số $40$ thỏa mãn $35 < 40 < 42$. Vậy có $40$ quả táo.
\textbf{Đáp án của câu hỏi là C}
}
\end{ex}

\begin{ex} %[2D1V01]
Hiệu giữa số lớn nhất có hai chữ số và số tròn chục nhỏ nhất có hai chữ số là
\choice
{\True $89$}
{$90$}
{$80$}
{$79$}
\loigiai{
Số lớn nhất có hai chữ số là $99$. Số tròn chục nhỏ nhất có hai chữ số là $10$.\nHiệu là: $99 - 10 = 89$.
\textbf{Đáp án của câu hỏi là A}
}
\end{ex}

\begin{ex} %[2D1V01]
Bạn Nam viết các số có hai chữ số giống nhau. Nam đã viết được bao nhiêu số?
\choice
{$10$ số}
{$8$ số}
{$11$ số}
{\True $9$ số}
\loigiai{
Các số có hai chữ số giống nhau là: $11, 22, 33, 44, 55, 66, 77, 88, 99$. Có tất cả $9$ số.
\textbf{Đáp án của câu hỏi là D}
}
\end{ex}

\begin{ex} %[2D1V01]
Trong các số sau: $48, 56, 84, 39, 93$, số có giá trị của chữ số hàng chục lớn nhất là
\choice
{$84$}
{\True $93$}
{$56$}
{$48$}
\loigiai{
Chữ số hàng chục của các số lần lượt là: 4, 5, 8, 3, 9. Chữ số lớn nhất là 9 (trong số 93).
\textbf{Đáp án của câu hỏi là B}
}
\end{ex}

\textbf{Phần 2. Trắc nghiệm đúng sai}
\setcounterpageref{ex}{0}

\begin{ex} %[2D1V01]
Xét tính đúng sai của các nhận định về cấu tạo số trong phạm vi 100:
\choiceTF
{\True Số bé nhất có hai chữ số khác nhau là 10.}
{Có tất cả 10 số tròn chục có hai chữ số.}
{\True Trong số 77, giá trị của chữ số 7 hàng chục hơn giá trị của chữ số 7 hàng đơn vị là 63 đơn vị.}
{\True Số 99 là số liền trước của số 100.}
\loigiai{
\begin{itemize}
    \item Ý a) Đúng: Vì chữ số hàng chục là 1, chữ số hàng đơn vị là 0 khác 1 nên 10 là số bé nhất có hai chữ số khác nhau.
    \item Ý b) Sai: Chỉ có 9 số tròn chục có hai chữ số từ 10 đến 90.
    \item Ý c) Đúng: Giá trị chữ số 7 hàng chục là 70, hàng đơn vị là 7. Ta có $70 - 7 = 63$.
    \item Ý d) Đúng: Ta có $100 - 1 = 99$, vậy 99 là số liền trước của 100.
\end{itemize}
\begin{itemize}
    \item a) \textbf{Đúng}.
    \item b) \textbf{Sai}.
    \item c) \textbf{Đúng}.
    \item d) \textbf{Đúng}.
\end{itemize}
}
\end{ex}

\begin{ex} %[2D1V01]
Cho các chữ số $3, 5, 8$. Lập các số có hai chữ số từ ba chữ số trên:
\choiceTF
{\True Lập được tất cả 6 số có hai chữ số khác nhau.}
{Số lớn nhất có thể lập được là số 85.}
{\True Tổng của số bé nhất và số lớn nhất trong các số có hai chữ số khác nhau lập được bằng 120.}
{Trong các số lập được có đúng 2 số chẵn.}
\loigiai{
\begin{itemize}
    \item Ý a) Đúng: Các số là $35, 38, 53, 58, 83, 85$ (tất cả 6 số).
    \item Ý b) Sai: Nếu cho phép chữ số giống nhau thì số lớn nhất là 88; nếu khác nhau là 85.
    \item Ý c) Đúng: Số bé nhất là 35, lớn nhất là 85: $35 + 85 = 120$.
    \item Ý d) Sai: Chỉ có 2 số có đuôi 8 là 38 và 58 là số chẵn.
\end{itemize}
\begin{itemize}
    \item a) \textbf{Đúng}.
    \item b) \textbf{Sai}.
    \item c) \textbf{Đúng}.
    \item d) \textbf{Sai}.
\end{itemize}
}
\end{ex}

\textbf{Phần 3. Trả lời ngắn}
\setcounterpageref{ex}{0}

\begin{ex} %[2D1V01]
Tìm số nhỏ nhất có hai chữ số mà tích hai chữ số của nó bằng 12.
\shortans{26}
\loigiai{
\textbf{Đáp số:} $26$\\
Các cặp số có tích bằng 12 là $2 \times 6 = 12$ và $3 \times 4 = 12$. Số nhỏ nhất là 26.
}
\end{ex}

\begin{ex} %[2D1V01]
Có bao nhiêu số có hai chữ số mà tổng hai chữ số của nó bằng 8?
\shortans{8}
\loigiai{
\textbf{Đáp số:} $8$\\
Các số đó là: 17, 26, 35, 44, 53, 62, 71, 80 (tất cả 8 số).
}
\end{ex}

\begin{ex} %[2D1V01]
Tìm số lớn nhất có hai chữ số khác nhau mà chữ số hàng chục gấp đôi chữ số hàng đơn vị.
\shortans{84}
\loigiai{
\textbf{Đáp số:} $84$\\
Hàng chục gấp đôi hàng đơn vị, để lớn nhất ta chọn hàng chục là 8, khi đó hàng đơn vị là $8 : 2 = 4$. Số đó là 84.
}
\end{ex}

\begin{ex} %[2D1V01]
Bạn An có một số que tính, nếu bó thành từng bó 10 que thì được 4 bó và còn thừa 7 que. Hỏi An có tất cả bao nhiêu que tính?
\shortans{47}
\loigiai{
\textbf{Đáp số:} $47$\\
$4$ bó là $40$ que. Thêm 7 que là $40 + 7 = 47$ que.
}
\end{ex}

\begin{ex} %[2D1V01]
Tìm một số biết rằng lấy số lớn nhất có hai chữ số trừ đi số đó thì được số tròn chục nhỏ nhất.
\shortans{89}
\loigiai{
\textbf{Đáp số:} $89$\\
Số lớn nhất có hai chữ số là 99. Số tròn chục nhỏ nhất là 10. Số đó là: $99 - 10 = 89$.
}
\end{ex}

\begin{ex} %[2D1V01]
Từ số 15 đến số 45 có tất cả bao nhiêu số tự nhiên liên tiếp?
\shortans{31}
\loigiai{
\textbf{Đáp số:} $31$\\
Số lượng số hạng: $(45 - 15) + 1 = 31$ số.
}
\end{ex}
